% !TeX root = ../../main.tex
% sections/app/appendix_cpi_derivation.tex

\chapter{総消費価格指数と財バスケットへの需要関数の導出（第2段階）}
\label{chap:appendix_cpi_derivation}

\section*{定理3：財バスケットへの需要関数と総消費価格指数}
\label{sec:appendix_cpi_derivation_theorem3}

家計\(h\)が、ある時点\(t\)において、所与の量の総消費バスケット\(c_t^{h\to W}\)を、国内財バスケットの真の価格指数\(p_t^H\)と外国財バスケットの真の価格指数\(p_t^{F*}\)の下で費用最小化的に購入するとき、以下の関係が成立する。
\begin{enumerate}
    \item\label{itm:appendix_cpi_derivation_theorem3_1}国内財バスケット\(c_t^{h\to H}\)および外国財バスケット\(c_t^{h\to F}\)への需要関数は、それぞれ次式で与えられる。
    \begin{equation}
    c_t^{h\to H}=\alpha^H\frac{p_t^{H\to W}}{p_t^H}c_t^{h\to W}
    \label{eq:appendix_cpi_derivation_theorem3_demand_domestic}
    \end{equation}
    \begin{equation}
    c_t^{h\to F}=(1-\alpha^H)\frac{p_t^{H\to W}}{e_tp_t^{F*}}c_t^{h\to W}
    \label{eq:appendix_cpi_derivation_theorem3_demand_foreign}
    \end{equation}
    \item\label{itm:appendix_cpi_derivation_theorem3_2}上記の需要関数に現れる\(p_t^{H\to W}\)は、総消費バスケットの価格指数（CPI）であり、各財バスケットの真の価格指数から以下のように定義される。
    \begin{equation}
    p_t^{H\to W}=(p_t^H)^{\alpha^H}(e_tp_t^{F*})^{1-\alpha^H}
    \label{eq:appendix_cpi_derivation_theorem3_cpi_def}
    \end{equation}
\end{enumerate}

\section*{証明}
\label{sec:appendix_cpi_derivation_proof}

この定理を、自国家計\(h\)の場合について証明する。

\subsection*{1.費用最小化問題の設定}
\label{sec:appendix_cpi_derivation_sub1}
この第2段階では、家計が時点\(t\)において、所与の総消費量\(c_t^{h\to W}\)を達成するために、国内財バスケット\(c_t^{h\to H}\)と海外財バスケット\(c_t^{h\to F}\)をどのように組み合わせれば総費用を最小化できるかを分析する。
\begin{itemize}
    \item\label{itm:appendix_cpi_derivation_sub1_obj}\textbf{最小化すべき総費用}:
    \begin{equation}
    \min_{c_t^{h\to H},c_t^{h\to F}}p_t^Hc_t^{h\to H}+e_tp_t^{F*}c_t^{h\to F}
    \label{eq:appendix_cpi_derivation_sub1_objective}
    \end{equation}
    \item\label{itm:appendix_cpi_derivation_sub1_const}\textbf{制約条件}（目標とする正規化されたバスケットの量）:
    \begin{equation}
    c_t^{h\to W}\equiv\frac{(c_t^{h\to H})^{\alpha^H}(c_t^{h\to F})^{1-\alpha^H}}{(\alpha^H)^{\alpha^H}(1-\alpha^H)^{1-\alpha^H}}
    \label{eq:appendix_cpi_derivation_sub1_constraint}
    \end{equation}
\end{itemize}

\subsection*{2.ラグランジュ関数と一階の条件（FOC）}
\label{sec:appendix_cpi_derivation_sub2}
制約式の対数を取ると扱いやすい。ラグランジュ関数\(\mathcal{L}^{h\to W}\)を設定する（この段階のラグランジュ乗数を\(\eta_t^h\)とする）。



\begin{equation}
\mathcal{L}^{h\to W}=p_t^Hc_t^{h\to H}+e_tp_t^{F*}c_t^{h\to F}-\eta_t^h\left(\alpha^H\ln c_t^{h\to H}+(1-\alpha^H)\ln c_t^{h\to F}-\ln c_t^{h\to W}-\text{const.}\right)
\label{eq:appendix_cpi_derivation_sub2_lagrangian}
\end{equation}
これを\(c_t^{h\to H}\)と\(c_t^{h\to F}\)でそれぞれ偏微分してゼロと置くと、以下の一階の条件（FOC）が得られる。
\begin{align}
\frac{\partial\mathcal{L}^{h\to W}}{\partial c_t^{h\to H}}=p_t^H-\eta_t^h\frac{\alpha^H}{c_t^{h\to H}}=0\quad&\implies\quad p_t^Hc_t^{h\to H}=\alpha^H\eta_t^h \label{eq:appendix_cpi_derivation_sub2_foc_h}\\
\frac{\partial\mathcal{L}^{h\to W}}{\partial c_t^{h\to F}}=e_tp_t^{F*}-\eta_t^h\frac{1-\alpha^H}{c_t^{h\to F}}=0\quad&\implies\quad e_tp_t^{F*}c_t^{h\to F}=(1-\alpha^H)\eta_t^h \label{eq:appendix_cpi_derivation_sub2_foc_f}
\end{align}

\subsection*{3.需要関数の導出（\(\eta_t^h\)を含む形）}
\label{sec:appendix_cpi_derivation_sub3}
上記の一階の条件をそれぞれ\(c_t^{h\to H}\)と\(c_t^{h\to F}\)について解くと、ラグランジュ乗数\(\eta_t^h\)を含む形で各バスケットへの需要関数が得られる。
\begin{align}
c_t^{h\to H}=\frac{\alpha^H\eta_t^h}{p_t^H}\quad,\quad c_t^{h\to F}=\frac{(1-\alpha^H)\eta_t^h}{e_tp_t^{F*}}
\label{eq:appendix_cpi_derivation_sub3_demand_with_multiplier}
\end{align}

\subsection*{4.ラグランジュ乗数\(\eta_t^h\)の導出}
\label{sec:appendix_cpi_derivation_sub4}
ステップ3で得た需要関数を、制約条件である総消費バスケットの定義式に代入する。
\begin{align}
c_t^{h\to W}&=\frac{1}{(\alpha^H)^{\alpha^H}(1-\alpha^H)^{1-\alpha^H}}\left(\frac{\alpha^H\eta_t^h}{p_t^H}\right)^{\alpha^H}\left(\frac{(1-\alpha^H)\eta_t^h}{e_tp_t^{F*}\!}\right)^{1-\alpha^H} \label{eq:appendix_cpi_derivation_sub4_multiplier_step1} \\
&=\frac{(\eta_t^h)^{\alpha^H+(1-\alpha^H)}}{(\alpha^H)^{\alpha^H}(1-\alpha^H)^{1-\alpha^H}}\cdot\frac{(\alpha^H)^{\alpha^H}(1-\alpha^H)^{1-\alpha^H}}{(p_t^H)^{\alpha^H}(e_tp_t^{F*})^{1-\alpha^H}} \label{eq:appendix_cpi_derivation_sub4_multiplier_step2} \\
c_t^{h\to W}&=\eta_t^h\cdot\frac{1}{(p_t^H)^{\alpha^H}(e_tp_t^{F*})^{1-\alpha^H}} \label{eq:appendix_cpi_derivation_sub4_multiplier_step3}
\end{align}
この式をラグランジュ乗数\(\eta_t^h\)について解く。
\begin{equation}
\eta_t^h=c_t^{h\to W}(p_t^H)^{\alpha^H}(e_tp_t^{F*})^{1-\alpha^H}
\label{eq:appendix_cpi_derivation_sub4_multiplier_result}
\end{equation}

\subsection*{5.総合物価指数\(p_t^{H\to W}\)と最終的な需要関数の導出}
\label{sec:appendix_cpi_derivation_sub5}
名目総消費は、最小化すべき総費用\(p_t^Hc_t^{h\to H}+e_tp_t^{F*}c_t^{h\to F}\)として定義される。ステップ2のFOCより、これは\(\alpha^H\eta_t^h+(1-\alpha^H)\eta_t^h=\eta_t^h\)に等しい。一方で、名目総消費は\(p_t^{H\to W}c_t^{h\to W}\)とも定義されるため、\(p_t^{H\to W}c_t^{h\to W}=\eta_t^h\)という関係が成立しなければならない。

この関係式にステップ4で導出した\(\eta_t^h\)の式を代入する。
\begin{equation}
p_t^{H\to W}c_t^{h\to W}=c_t^{h\to W}(p_t^H)^{\alpha^H}(e_tp_t^{F*})^{1-\alpha^H}
\label{eq:appendix_cpi_derivation_sub5_equivalence}
\end{equation}
両辺の\(c_t^{h\to W}\)を消去することで、定理の項目2で示された\textbf{総合物価指数\(p_t^{H\to W}\)}が導かれる。
\begin{equation}
p_t^{H\to W}=(p_t^H)^{\alpha^H}(e_tp_t^{F*})^{1-\alpha^H}
\label{eq:appendix_cpi_derivation_sub5_cpi_final}
\end{equation}
最後に、この\(\eta_t^h=p_t^{H\to W}c_t^{h\to W}\)という関係をステップ3で得た需要関数に代入することで、定理の項目1で示された最終的な財バスケットへの需要関数が完成する。
\begin{equation}
c_t^{h\to H}=\frac{\alpha^H(p_t^{H\to W}c_t^{h\to W})}{p_t^H}=\alpha^H\frac{p_t^{H\to W}}{p_t^H}c_t^{h\to W}
\label{eq:appendix_cpi_derivation_sub5_demand_domestic_final}
\end{equation}
\begin{equation}
c_t^{h\to F}=\frac{(1-\alpha^H)(p_t^{H\to W}c_t^{h\to W})}{e_tp_t^{F*}}=(1-\alpha^H)\frac{p_t^{H\to W}}{e_tp_t^{F*}}c_t^{h\to W}
\label{eq:appendix_cpi_derivation_sub5_demand_foreign_final}
\end{equation}
（証明終）